\documentclass[11pt]{article}
\usepackage{ap-psych}
\title{Unit 5: Mental and Physical Health: Question Set}
\author{Paul Wang}
\date{November 2026}
\docpart{Unit 5 \textperiodcentered\ Question Set}

\begin{document}
\maketitle

\begin{abstract}
Exam-style multiple choice on all five Unit 5 topics, with a shared-study set on a treatment trial and a full Evidence-Based Question on exercise and depression. Disorder items give symptom descriptions and ask for the best-fitting category; treatment items describe what a therapist does and ask for the approach or mechanism. All items and sources are original and fictional, written for practice, and the vignettes are for exam study, not diagnosis. Allow about 30 minutes for the multiple choice and 45 minutes for the EBQ.
\end{abstract}

\scoringline{MCQ: \blank\ / \qtotal \qquad EBQ: \blank\ / 7 \qquad Target: 21+ MCQ and 5+ EBQ \qquad Pace: 72 s per MCQ}

\section{Stress, coping, and positive psychology (5.1--5.2)}

\begin{q}{2}{5.1 GAS}{Immediate surge of sympathetic arousal = alarm stage.}
Seconds after hearing a fire alarm, Pedro's heart pounds and his breathing quickens. In Selye's model, he is in the
\chh{resistance stage}{alarm stage}{exhaustion stage}{recovery stage}
\end{q}

\begin{q}{4}{5.1 Coping}{Managing feelings rather than the stressor = emotion-focused coping.}
Unable to change her company's layoffs, Gina meditates and talks with friends to manage her worry. This is
\ch{problem-focused coping}{the tend-and-befriend response}{a catastrophe}{emotion-focused coping}
\end{q}

\begin{q}{1}{5.1 Stress and health}{Chronic stress hormones suppress immune function.}
Students during final exams take longer to heal from minor wounds than during vacation. The best explanation is that
\ch{prolonged stress suppresses immune system functioning}{eustress speeds healing}{exams activate the parasympathetic nervous system}{daily hassles have no physical effects}
\end{q}

\begin{q}{3}{5.1 Tend-and-befriend}{Responding to stress by nurturing and seeking support.}
After a frightening earthquake, many residents respond by checking on neighbors and gathering with family. This pattern is called
\chh{fight-or-flight}{the exhaustion stage}{tend-and-befriend}{learned helplessness}
\end{q}

\begin{q}{2}{5.2 Positive psychology}{Gratitude practices raise subjective well-being.}
A study finds that students who wrote weekly gratitude letters reported greater life satisfaction than controls. This finding is most closely associated with
\chh{psychodynamic theory}{positive psychology}{the medical model}{behavioral perspective on disorders}
\end{q}

\section{Defining and explaining disorders (5.3)}

\begin{q}{4}{5.3 Definitions}{Behavior typical within its cultural context is not deviant.}
Hearing the voice of a deceased relative during a culturally accepted mourning ritual would most likely not be considered a disorder because it is
\chh{distressful}{dysfunctional}{dangerous}{not deviant within its cultural context}
\end{q}

\begin{q}{1}{5.3 Perspectives}{Maladaptive thinking patterns as the cause = cognitive perspective.}
A psychologist attributes a client's depression to her habit of interpreting every setback as proof she is worthless. This reflects the
\chh{cognitive perspective}{sociocultural perspective}{biological perspective}{humanistic perspective}
\end{q}

\begin{q}{3}{5.3 Labels}{Labels guide treatment and research but can bias perception and create stigma.}
A main criticism of diagnostic labels is that they
\ch{prevent researchers from studying disorders}{are never used internationally}{can bias how others perceive and treat the person}{have no use in choosing treatment}
\end{q}

\section{Categories of disorders (5.4)}

\begin{q}{2}{5.4 Panic disorder}{Sudden unexpected attacks + worry about more = panic disorder.}
Without warning, Erin suddenly feels her heart race, can't breathe, and is sure she is dying; this has happened several times, and she now worries constantly about the next attack. She most likely has
\chh{specific phobia}{panic disorder}{PTSD}{bipolar I}
\end{q}

\begin{q}{4}{5.4 GAD}{Chronic, diffuse worry about many things = generalized anxiety disorder.}
For over a year, Sam has worried uncontrollably about money, health, school, and his family, even when things are going well. The best fit is
\ch{social anxiety disorder}{obsessive-compulsive disorder}{agoraphobia}{generalized anxiety disorder}
\end{q}

\begin{q}{1}{5.4 PTSD}{Trauma, flashbacks, avoidance, hypervigilance > 1 month = PTSD.}
Months after a combat deployment, a veteran has nightmares and flashbacks, avoids crowds, and startles at loud noises. These symptoms most closely match
\chh{PTSD}{schizophrenia}{dissociative identity disorder}{ADHD}
\end{q}

\begin{q}{3}{5.4 Schizophrenia}{Absence of normal emotion and speech = negative symptoms.}
A patient speaks very little, shows almost no facial expression, and has withdrawn from friends. These are best described as
\ch{positive symptoms of schizophrenia}{signs of a manic episode}{negative symptoms of schizophrenia}{symptoms of a specific phobia}
\end{q}

\begin{q}{2}{5.4 Dissociative}{Memory loss for personal identity with unplanned travel = dissociative amnesia with fugue.}
A man is found in another state with no memory of who he is or how he got there, following a traumatic event. This is most consistent with
\ch{schizophrenia}{dissociative amnesia with fugue}{bipolar II}{retrograde amnesia from a concussion}
\end{q}

\begin{q}{4}{5.4 Personality}{Unstable relationships, self-image, mood, fear of abandonment = borderline.}
Hana's relationships swing between idealizing and devaluing partners; her self-image and moods shift rapidly, and she fears being abandoned. This pattern is most consistent with
\ch{schizoid personality disorder}{antisocial personality disorder}{avoidant personality disorder}{borderline personality disorder}
\end{q}

\begin{q}{1}{5.4 Neurodevelopmental}{Childhood-onset social-communication deficits + restricted, repetitive behaviors = ASD.}
A 4-year-old rarely makes eye contact, has difficulty with back-and-forth conversation, and becomes very upset if his toys are moved out of a precise line. These features are most consistent with
\ch{autism spectrum disorder}{ADHD}{social anxiety disorder}{schizotypal personality disorder}
\end{q}

\begin{q}{3}{5.4 Biology}{Schizophrenia linked to excess dopamine activity; antipsychotics block it.}
The dopamine hypothesis is supported by the finding that
\ch{SSRIs reduce hallucinations}{lithium cures schizophrenia}{drugs that block dopamine receptors reduce positive symptoms}{low GABA causes delusions}
\end{q}

\section{Treatment (5.5)}

\begin{stimulus}
\textbf{Study F.} 180 adults diagnosed with major depressive disorder were randomly assigned to 16 weeks of CBT, an SSRI, or a pill placebo with clinical check-ins. Clinicians who rated depression did not know participants' conditions. Mean drop in depression score (higher = more improvement): CBT, 11.2; SSRI, 10.8; placebo, 5.1. CBT and SSRI did not differ significantly; both differed from placebo (\(p < .01\)). One year after treatment ended, relapse rates were 24\% (CBT) and 47\% (SSRI, medication stopped).
\end{stimulus}

\begin{q}{2}{5.5 Data}{CBT ≈ SSRI at 16 weeks; both > placebo.}
Which conclusion about the 16-week results of Study F is best supported?
\ch{The SSRI was significantly more effective than CBT.}{CBT and the SSRI were similarly effective, and both outperformed the placebo.}{The placebo had no effect.}{CBT was ineffective.}
\end{q}

\begin{q}{4}{5.5 Data}{Lower relapse after CBT suggests skills learned in therapy persist.}
The relapse data in Study F are most consistent with the idea that
\ch{medication changes thinking patterns permanently}{placebo effects last longest}{SSRIs work only through expectations}{CBT teaches skills that continue to protect after treatment ends}
\end{q}

\begin{q}{1}{5.5 Methods}{Raters blind to condition control experimenter bias.}
In Study F, keeping the raters unaware of each participant's condition was intended to reduce
\chh{experimenter bias}{random assignment}{attrition}{cohort effects}
\end{q}

\begin{q}{3}{5.5 Behavioral}{Immediate, intense exposure without gradual steps = flooding.}
To treat a fear of elevators, a therapist has the client ride an elevator repeatedly for an hour on the first day, with no gradual buildup. This is
\chh{systematic desensitization}{aversive conditioning}{flooding}{free association}
\end{q}

\begin{q}{2}{5.5 Cognitive}{Disputing irrational ``musts'' = Ellis's REBT.}
A therapist forcefully challenges a client's belief that ``I must be liked by everyone or I am worthless.'' This technique is central to
\ch{person-centered therapy}{rational emotive behavior therapy}{psychoanalysis}{a token economy}
\end{q}

\begin{q}{4}{5.5 Biomedical}{Long-term antipsychotic use can cause involuntary movements (tardive dyskinesia).}
After years on an older antipsychotic drug, a patient develops involuntary lip-smacking and facial twitches. This side effect is
\chh{mania}{withdrawal}{catatonia}{tardive dyskinesia}
\end{q}

\begin{q}{1}{5.5 Biomedical}{ECT is used for severe, treatment-resistant depression.}
For a patient with severe depression who has not improved after several medications and therapies, a psychiatrist is most likely to consider
\ch{electroconvulsive therapy}{a lobotomy}{benzodiazepines alone}{flooding}
\end{q}

\section{Evidence-Based Question}

\begin{frq}{Question 2 \textperiodcentered\ EBQ \textperiodcentered\ 7 points \textperiodcentered\ 45 minutes}
\small
\textbf{Directions.} Read the three summarized sources (fictional, written for practice) and answer the prompt. Cite sources by author and year.

\textbf{Prompt.} Using evidence from at least two sources, develop and justify an argument about whether physical exercise is an effective way to reduce symptoms of depression.

\medskip
\textbf{Source 1.} Moreau \& Baptiste (2023). \emph{Participants}: 96 adults with mild-to-moderate depression in Belgium. \emph{Method}: Experiment. Participants were randomly assigned to supervised aerobic exercise (40 minutes, 3 times a week, 12 weeks) or a waitlist. \emph{Results}: Depression scores fell an average of 9.4 points in the exercise group and 2.1 points in the waitlist group (\(p < .01\)). Exercisers also reported more weekly pleasant activities.

\textbf{Source 2.} Adeyemi, Sato, \& Ruiz (2024). \emph{Method}: Meta-analysis of 38 randomized trials (total \(N = 3{,}410\)). \emph{Results}: Exercise produced a moderate reduction in depressive symptoms compared with control conditions. Effects were larger for supervised programs and similar in size to those of CBT in trials that compared them directly.

\textbf{Source 3.} Halvorsen (2022). \emph{Participants}: 6{,}200 Norwegian adults without depression at the start. \emph{Method}: Longitudinal correlational study over 8 years. \emph{Results}: Adults who reported at least 150 minutes of moderate activity per week were 26\% less likely to develop depression than inactive adults, after controlling for age, income, and health.

\medskip
\begin{enumerate}[label=(\Alph*),leftmargin=2em,itemsep=3pt]
\item Propose a specific and defensible claim based in psychological science that responds to the prompt.
\item (i) Support your claim using at least one piece of specific and relevant evidence from one of the sources. (ii) Explain how the evidence from Part B(i) supports your claim using a psychological perspective, theory, concept, or research finding learned in AP Psychology.
\item (i) Support your claim using an additional piece of specific and relevant evidence from a different source. (ii) Explain how the evidence from Part C(i) supports your claim using a different psychological perspective, theory, concept, or research finding learned in AP Psychology.
\end{enumerate}
\end{frq}

\makekey

\subsection*{EBQ scoring guide (practice rubric, 7 points)}
\small
Modeled on the published EBQ structure (claim, two pieces of cited evidence, reasoning with course content).
\begin{compare}{Part}{Earns credit}{Common miss}
A. Claim (1) & ``Regular supervised aerobic exercise reduces symptoms of mild-to-moderate depression in adults.'' & ``Exercise is healthy''; ``exercise cures all depression'' (overclaims the sources). \\
B(i). Evidence (1) & Moreau \& Baptiste (2023): randomly assigned exercisers' depression scores fell 9.4 points vs.\ 2.1 for the waitlist (\(p < .01\)). & ``A study found exercise helps'' with no source or result. \\
B(ii). Reasoning (2) & 1 pt: links evidence to claim; random assignment supports a causal conclusion. 1 pt: course concept, e.g., exercise increases \emph{endorphin}, \emph{serotonin}, and \emph{norepinephrine} activity (biological perspective), or more pleasant activities provide \emph{positive reinforcement} that counters withdrawal (behavioral perspective). & Listing neurotransmitters without connecting them to the score drop. \\
C(i). Evidence (1) & Adeyemi et al.\ (2024): across 38 randomized trials, exercise produced a moderate reduction in symptoms, comparable to CBT. \emph{Or} Halvorsen (2022): active adults were 26\% less likely to develop depression over 8 years. & A second detail from the same source used in B. \\
C(ii). Reasoning (2) & 1 pt: links to claim. 1 pt: different concept, e.g., a \emph{meta-analysis} pools many studies, making the conclusion more reliable; exercise as \emph{problem-focused/emotion-focused coping} that buffers stress (diathesis--stress model). If using Halvorsen, notes it is correlational, so it supports prevention but cannot prove causation. & Treating Halvorsen's correlation as proof that exercise prevents depression. \\
\end{compare}

\end{document}
